\section{Related Work}

Autoregressive models, exemplified by large language models, generate sequences token-by-token in a monotonic manner, i.e., once a token is emitted, later tokens must condition on it, and the model cannot revise earlier decisions~\cite {bachmann2024pitfalls, thoppilan2022lamda}.
A fundamental limitation of this paradigm is error propagation~\cite{zhang2023snowball, gan2025rethinking}, where a mistake made early in generation can propagate and amplify over the remainder of the sequence, since there is no mechanism for erasing flawed intermediate content.

\subsection{Post-hoc Verification and Iterative Refinement}
A common direction to tackle the error propagation issue is post-hoc self-critique and revision~\cite{shinn2023reflexion, welleck2022generating, gou2023critic, wang2022self}. These methods typically follow an iterative loop, i.e., generate an initial draft, produce feedback through the same model or an auxiliary model, and then regenerate a revised answer conditioned on the critique.
For instance, Self-Refine~\cite{madaan2023selfrefine} demonstrates that models can improve responses via multiple rounds of self-generated feedback, without task-specific fine-tuning or additional supervision: the same pretrained model alternates between drafting, critiquing, and revising. 
Reflexion~\cite{shinn2023reflexion} extends this idea to agentic settings by maintaining an explicit memory of past failures. 
The agent reflects on these errors and uses them as "lessons" to guide future attempts.
More recently, intrinsic self-verification aims to endow the model with an internal judgment of whether its partial reasoning or output is trustworthy, reducing reliance on external critics. 
ReVISE~\cite{lee2025revise} trains a model to decide whether a reasoning trajectory is sound or should be revised. 
While effective, post-hoc refinement and self-verification pipelines typically introduce additional inference overhead because they require at least one extra generation pass to critique and revise.

\subsection{Self-Refine via Backtrack Generation}
The most relevant line of work to our proposed method is self-refinement via backtrack generation, which allows the model to interleave generation with on-the-fly self-assessment and revision within a single decoding pass.
These works mainly focus on safety violations in autoregressive sequence models.
For example, \citet{zhang2025backtracking} proposes a backtracking-enabled decoding mechanism for safer generation by introducing a special \texttt{[RESET]} token and fine-tuning the model to emit it when the output drifts toward unsafe content. 
\citet{selreinforcement} extends this idea by introducing category tokens \texttt{[CATEGORY]} that flag the type of safety violation, enabling targeted revisions.
\citet{wu2025generate} introduces REVERSE for vision-language generation, training a model to interleave generation with self-refinement by inserting special tokens, i.e., \texttt{<UN>} and \texttt{<CN>}, that trigger verification and correction behaviors. 
Unlike prior backtracking methods that focus primarily on safety violations~\cite{zhang2025backtracking,selreinforcement} or vision-language tasks~\cite{wu2025generate}, \modelname targets complex reasoning domains where error propagation is particularly detrimental. 
Moreover, N-MARS employs a fine-grained non-monotonic generation strategy that allows the model to revise arbitrary segments of the generated sequence, rather than only from fixed positions, i.e., end of sentences or end of generation.
This flexibility is crucial for reasoning tasks, where errors can occur at any step, and earlier mistakes may require revisiting multiple subsequent tokens.



% Beyond post-hoc revision, another line of work seeks to mitigate error accumulation by making generation itself non-monotonic—allowing the model to retract or overwrite earlier content during decoding. If the model can detect a failure mode mid-generation and backtrack immediately, it can correct errors before they contaminate the remainder of the sequence.





